Для оценки точности энкодера в данной работе была проведена серия измерений при равномерном вращении кодового диска прецизионным шаговым двигателем по и против часовой стрелки. Эксперименты выполнялись как для каждой из двух диаметрально противоположных считывающих головок по отдельности, так и для их совместного использования. Для каждого отсчета определялись среднее значение абсолютного угла и среднеквадратическое отклонение по массиву всех кодовых слов, попавших в поле зрения матрицы.

Предварительно для каждой из считывающих головок была выполнена независимая калибровка и настройка. В рамках этой процедуры измерялись плоские поля, характеризующие виньетирование оптической системы, и оценивались коэффициенты дисторсии. Кроме того, по изображению в программной модели визуализации определялось эффективное фокусное расстояние для каждой головки ручным подбором по оптимальной и наиболее четкой фокусировке.

Анализ среднего значения угла позволяет оценить систематическую составляющую погрешности энкодера. Анализ среднеквадратического отклонения характеризует повторяемость и стабильность измерений. Сопоставление результатов, полученных при вращении в противоположных направлениях, дополнительно позволяет выявить возможную направленную асимметрию погрешности.

Усреднение результатов выполнялось по отдельным полным периодам сигнала. После выделения активного участка записи оценивался период, затем по линейному тренду определялись границы полных периодов, внутри которых сигнал приводился к общей фазовой сетке. Для каждой точки фазы вычислялась средняя детрендированная ошибка по всем периодам, что позволяло получить усредненную фазовую зависимость систематической составляющей погрешности.

Статистическая составляющая среднеквадратического отклонения определялась как выборочное СКО между отдельными периодами в каждой точке фазы. Инструментальная составляющая, соответствующая собственной погрешности алгоритма обработки для одной головки, вычислялась как среднеквадратическое значение СКО отдельных измерений, приведенных к той же фазовой сетке. Итоговое СКО для одной головки определялось как квадратичная сумма статистической и инструментальной составляющих:
\begin{equation*}
\sigma_{\mathrm{full}}(\varphi)=\sqrt{\sigma_{\mathrm{stat}}^2(\varphi)+\sigma_{\mathrm{instr}}^2(\varphi)}.
\end{equation*}

По результатам экспериментов было установлено, что среднеквадратическое отклонение зависит от направления измерений. Предположительно, наблюдаемая асимметрия связана с тем, что во всех экспериментах считывание ПЗС-матрицы выполнялось в одном и том же направлении. При совпадении направления считывания с направлением вращения диска значение СКО уменьшалось, тогда как при противоположных направлениях вращения и считывания оно возрастало.

Отдельно был исследован эффект усреднения измерений, используемый в самокалибровке методом путевого усреднения \ref{par:avg_path_method}. Поскольку рассматриваемая величина является угловой, для объединения оценок использовалось круговое среднее. По результатам такого объединения наблюдалось существенное подавление периодической систематической погрешности: наблюдалось существенное ослабление первой гармоники и в остаточной зависимости наиболее выраженной оставалась вторая гармоника (см. \hyperref[sec:application4]{Приложение 4}).

Для оценки разброса объединенного измерения использовалась величина
\begin{equation*}
\sigma_{\mathrm{avg}} = \frac{1}{2}\sqrt{\sigma_1^2 + \sigma_2^2}.
\end{equation*}
При предположении о некоррелированности ошибок двух головок и при их близких уровнях это приводит к уменьшению среднеквадратического отклонения результирующей оценки примерно в $\sqrt{2}$ раза, что было подтверждено измерениями -- уровень СКО снизился с $\sim0.08$ до $\sim0.06$ градусов.

По результатам обработки строились диагностические графики, отражающие усредненную периодическую ошибку, структуру составляющих погрешности и разброс отдельных реализаций относительно средней зависимости:

\begin{enumerate}
    \item детрендированная усредненная зависимость ошибки от фазы с границами среднеквадратического отклонения;
    \item фазовые зависимости статистической, инструментальной и полной составляющих погрешности;
    \item отдельные детрендированные реализации, наложенные на общую фазовую ось;
    \item сопоставление результатов, полученных для отдельных головок и для их кругового среднего.
\end{enumerate}

Максимальное отклонение измеряемого угла (точность), полученная усреднением показаний двух головок достигла в среднем не более $0.08$ градусов при вращении противоположно считыванию, воспроизводимость угла при вращении в обе стороны составляет $0.1$ градуса.

Все диагностические графики см. в \hyperref[sec:application4]{Приложении 4}.